% Original synthetic teaching sample for TeXada; AGPL-3.0-only.
% No copied textbook content or real personal data.
\documentclass[12pt]{article}
\usepackage{amsmath,amssymb}
\usepackage[margin=2.2cm]{geometry}
\setlength{\parindent}{0pt}
\setlength{\parskip}{0.7em}
% Intentional mathematical error with valid formula syntax.
% Teaching purpose: syntax parsing does not evaluate the truth of an equality.
\title{Calculus Notes: Syntax Is Not Truth}
\author{LinguistsWantTech --- synthetic teaching sample}
\date{}
\begin{document}
\maketitle
\section*{1. The deliberately false claim}
The following displayed-sized inline formula intentionally states
an incorrect integral value:
$\displaystyle \int_0^1 x\,dx=1$.
\section*{2. An independent mathematical check}
The antiderivative is $\frac{x^2}{2}$. Evaluating it at one and zero
gives $\frac{1}{2}$, not one.
\section*{3. What the software can establish}
The formula can have valid syntax while its equality is false.
The current syntax checker does not certify the mathematical statement
and should not invent a repair for a formula that it accepts.
\textbf{Teaching note.} This original negative control must remain
available for human review; zero syntax problems is not a correctness score.
\end{document}
